\labeledsection{Term Languages and Abstract Grammars}{sec:term_langs_abstract_grams}
\commandnote{
In most applications of \linkterm{symbolic AI}{symbolicai}, the \linkterm{formal languages}{formal_language} are very strucutred.
}

\textbf{Example (\linkterm{Arithmetic}{arithmetic} Expressions): }Consider the \linkterm{CFG}{contex_free_grammar} $G$ and the $G$-derivation:

\begin{verbatim}
term   → num | var | sum | prod | neg
num    → digit*
digit  → 1 | 2 | 3 | 4 | 5 | 6 | 7 | 8 | 9 | 0
var    → 'x' num
sum    → '(' term '+' term ')'
prod   → '(' term '*' term ')'
neg    → '-(' term ')'

term → G -(term)
     → G -((term * term))
     → G -(((term + term) * num))
     → G -(((var + num) * digit digit))
     → G -(((x num + digit*) * 555))
     → G -(((x digit digit + 3) * 555))
     → G -(((x29 + 3) * 555))
\end{verbatim}
$G$ accepts the string $-(((X29+3) * 555))$ but not $-(X29*3($.

\commandnote{
We call such \linkterm{languages}{formal_language} term languages
}

\textbf{Observation: }Strings in $L(G)$ are \textit{well-bracketed} and can be split into sub-strings from $L(G)$.

\textbf{Intuition: }The parse \linkterm{trees}{tree} are the primary \linkterm{objects}{math_object} for \linkterm{symbolic AI}{symbolicai}, the strings are just the technical I/O.

\definition{Parse Tree}{
Let $G$ be a \linkterm{CFG}{contex_free_grammar} that accepts a string $s$. Then a \textbf{parse tree} is a \linkterm{node labeled}{labeled_graph}, ordered \linkterm{tree}{tree} $P_s$ that represents the syntactic strucutre of $s$.
}{parse_tree}
\textbf{Example: }the \linkterm{parse tree}{parse_tree} of "the teacher sleeps":
\begin{verbatim}
        S
       / \
     NP   Vi
    /  \    \
Article N   sleeps
   |    |
  the teacher  
\end{verbatim}

\commandnote{
There may be multiple derivations that accept a string $s$ in a \linkterm{CFG}{contex_free_grammar}. For example, $S \to \text{ NP } \text{Vi} \to \text{ Article N } \text{Vi} \cdots$ or $S \to \text{ NP } \text{Vi} \to \text{ NP } \text{sleeps} \cdots$, and so on.
}

Let $G$ be \linkterm{CFG}{contex_free_grammar} and $D := s \to_{G}^{*} t$ a $G$-derivation where $t$ is a \linkterm{sentence}{sentence_grammar} of $G$. We define the \linkterm{parse tree}{parse_tree} $P_D$ of $D$ \linkterm{recursively}{recursion}:
\begin{itemize}
\item IF $D = s \to_{G}^{p} t$, then $p$ is a \linkterm{lexical}{lexical_grammar} \linkterm{rule}{phrase_structure_grammar} $s\to t$ and $t$ consists of a single \linkterm{terminal}{phrase_structure_grammar}. Then $P_D$ is the \linkterm{tree}{tree} whose root has label $s$ with one child labeled with $t$.
\item If $D = s \to_{G}^{p} \bar{s} \to_{G}^{*} t$, $\bar{s} \to_{G}^{*} t$, and $p = h \to a_1 \cdots a_n$, then the root of $P_D$ has label $h$ and it has children that are the \linkterm{parse trees}{parse_tree} for the subderivations for the $a_i$ in $\bar{D}$.
\end{itemize}

\commandnote{Term languages are ones that have enough structural characters so that the \linkterm{parse tree}{parse_tree} is obvious.}

In \linkterm{SML}{def:sml} we can program the same \linkterm{grammar}{phrase_structure_grammar} as an abstract datatype:
\begin{verbatim}
datatype aterm =
    anum of int
  | avar of int
  | aneg of aterm
  | asum of aterm * aterm
  | aprod of aterm * aterm;
\end{verbatim}

And we can express arithmetic expressions as \linkterm{SML}{def:sml} expressions directly, e.g. $-(((X29+3) * 555))$:
\begin{verbatim}
- val ex = aneg (aprod (asum (avar 29, anum 3), anum 555));
val ex = aneg (aprod (asum (avar 29,anum 3),anum 555)) : aterm
\end{verbatim}

Moreover, we can traverse the term \linkterm{recursively}{recursion} and print it as a string:
\begin{verbatim}
fun aprint (anum n) = Int.toString n
  | aprint (avar n) = "x" ^ Int.toString n
  | aprint (aneg t) = "-(" ^ aprint t ^ ")"
  | aprint (asum (t1, t2)) = "(" ^ aprint t1 ^ "+" ^ aprint t2 ^ ")"
  | aprint (aprod (t1, t2)) = "(" ^ aprint t1 ^ "*" ^ aprint t2 ^ ")";
\end{verbatim}

Let's test that on our \texttt{ex}:
\begin{verbatim}
- aprint ex;
val it = "-(((x29+3)*555))" : string
\end{verbatim}

We can also do more complex tasks via \linkterm{tree}{tree} traversal. For example, we can evaluate \texttt{"-(((x29+3)*555))"}. To do that, we need to know the value of the variable $X29$. This is traditionally given by an assignment that assigns values to variables, e.g. $X1 \mapsto 3, X29 \mapsto 5, \cdots$ (a dictionary), which we can represent as a list of pairs in \linkterm{SML}{def:sml}:
\begin{verbatim}
type assignment = (int * aterm) list
\end{verbatim}
Then we can represent the evaluation function by \linkterm{recursion}{recursion} over the structure of the expression:
\begin{verbatim}
fun aeval (anum n, a:assignment) = n
|   aeval (aneg (n), a) = ~(aeval (n,a))
|   aeval (asum (s,t), a) = aeval (s,a) + aeval (t,a)
|   aeval (aprod (s,t), a) = aeval (s,a) * aeval (t,a);
\end{verbatim}
In all cases, the assignment is just passed on to the \linkterm{recursive call}{recursion}. The only exception is the variable case, where we need to \linkterm{recursively}{recursion} find the right value from the assignment, so the final function would look like:
\begin{verbatim}
fun aeval (anum n, a) = n
  | aeval (aneg t, a) = ~ (aeval (t, a))
  | aeval (asum (s, t), a) = aeval (s, a) + aeval (t, a)
  | aeval (aprod (s, t), a) = aeval (s, a) * aeval (t, a)
  | aeval (avar n, []) = raise Fail "unbound variable"
  | aeval (avar n, (m, r) :: l) =
      if n = m then aeval (r, l) else aeval (avar n, l);
\end{verbatim}

Let's first assign $X29 \mapsto 5$
\begin{verbatim}
val env = [(29, anum 5)];
\end{verbatim}
Then test the function on our term:
\begin{verbatim}
- val ex = aneg (aprod (asum (avar 29, anum 3), anum 555));
val ex = aneg (aprod (asum (avar 29,anum 3),anum 555)) : aterm

- val result = aeval (ex, env);
val result = ~4440 : int
\end{verbatim}

\commandnote{
The situation for \linkterm{arithmetic}{arithmetic} expressions in \linkterm{SML}{def:sml} is typical for term languages.
\begin{itemize}
    \item Inductive data types, one constructor per \linkterm{production rule}{phrase_structure_grammar}
    \item Expressions as \linkterm{tree}{tree}-structured data structures (no brackets needed)
    \item All computation on expressions via \linkterm{recursive}{recursion} \linkterm{tree}{tree} traversal
\end{itemize}
}

\definition{Tree Grammar}{
A tree grammar is a \linkterm{grammar}{phrase_structure_grammar} where the result data type is \linkterm{trees}{tree} not strings.
}{tree_grammar}
If we interpret the brackets as "\linkterm{parse tree}{parse_tree} constructors", then we do not have to worry about matching them. The string representation - with e.g. infix notations for operators - is just a derived / secondary artifact for storage / communication (implementation details).

\definition{Abstract vs Concrete Grammars}{
The underlying \linkterm{tree}{tree} \linkterm{grammar}{phrase_structure_grammar} of a \linkterm{language}{formal_language} is often called the \textbf{abstract grammar}; it captures the essence of the \linkterm{formal language}{formal_language} (the constructors), and any derived (string) \linkterm{grammar}{phrase_structure_grammar} - there may be multiple - is called a \textbf{concrete grammar}; they correspond to, e.g. the results of print functions.
}{abstract_concrete_grammar}

\commandnote{We are mostly interested in \linkterm{\textbf{abstract grammars}}{abstract_concrete_grammar} in \linkterm{symbolic AI}{symbolicai}}

